%% kapa: exam=1 round=33 year=2022 period=2 form=A subject=accounting subject_ko=회계학 seq=15 number=55
%% source: Q-Net 감정평가사 기출문제 (cst00309) artlSeq=5209997 file=raw/1008/5209997/2022년도 제33회 감정평가사 1차시험 2교시(감정평가 관계법규, 회계학).pdf
%% conversion: kordoc 4.24.51 -> tools/gen1.py (auto); review=proofread-sonnet; flags=none
\begin{question}{55}
(주)감평은 20×1년 10월 1일에 고객과 원가 ￦900의 제품을 ￦1,200에 판매하는 계약을 체결하고 즉시 현금 판매하였다. 계약에 따르면 (주)감평은 20×2년 3월 31일에 동 제품을 ￦1,300에 재매입할 수 있는 콜옵션을 보유하고 있다. 동 거래가 다음의 각 상황에서 (주)감평의 20×2년도 당기순이익에 미치는 영향은? (단, 각 상황(A, B)은 독립적이고, 화폐의 시간가치는 고려하지 않으며, 이자비용(수익)은 월할계산한다.)
\par{\small\centering\begin{tabular}{|c|l|}\hline
상황 & \multicolumn{1}{c|}{내 용} \\ \hline
A & 20×2년 3월 31일에 (주)감평이 계약에 포함된 콜옵션을 행사한 경우 \\
B & 20×2년 3월 31일에 계약에 포함된 콜옵션이 행사되지 않은 채 소멸된 경우 \\ \hline
\end{tabular}\par}
\begin{choices}
  \item 상황A ￦100 감소 \quad 상황B ￦100 증가
  \item 상황A ￦50 감소 \quad 상황B ￦100 증가
  \item 상황A ￦50 감소 \quad 상황B ￦350 증가
  \item 상황A ￦300 증가 \quad 상황B ￦350 증가
  \item 상황A ￦400 증가 \quad 상황B ￦400 증가
\end{choices}
\end{question}
\answer{3}
